\section*{Disponibilidade de artefatos}
O repositório está em \url{https://github.com/marcelocorni/anomalias-ae-if-2025}. O código e as dependências do experimento estão fixados no \href{https://github.com/marcelocorni/anomalias-ae-if-2025/tree/852f5452d4f59b2e4615914b69ce0b7a19d423ef}{commit \texttt{852f545}}. O perfil \path{configuracao/execucao_revisao_01.json} aponta para a configuração que registra o SHA-256 do CSV histórico. As entradas externas são os Parquets da AWS e o CSV exato identificado por esse checksum. A \href{https://github.com/marcelocorni/anomalias-ae-if-2025/blob/852f5452d4f59b2e4615914b69ce0b7a19d423ef/documentacao/execucao-etapa4.md}{documentação de execução} descreve a preparação e o comando do perfil.

Os \href{https://github.com/marcelocorni/anomalias-ae-if-2025/tree/9a6df232b9141f32cc86a5ccb0bd50a459e7c2c8/artigo}{materiais publicados com a revisão anterior} incluem \path{artigo/dados-artigo.json} e \path{artigo/suplemento/orcamentos.csv}. Esses agregados registram taxas, interseções, PSI e os valores de $k$, P@k, recall e lift. O suplemento \path{suplemento/concentracao-remetentes.json} registra a concentração, a assinatura da execução e o hash da referência normalizada. A reprodução das curvas Precision-Recall usa também os pontos dos relatórios de avaliação \path{labelcloud/}, gerados pelo perfil experimental.
